\section{Introduction}\label{sec:introduction}

The H\'enon--Lane--Emden system is
\begin{equation}\label{eq:system}
 -\Delta u=|x|^A v^p,\qquad
 -\Delta v=|x|^B u^q
 \qquad\text{in }\R^n\setminus\{0\}.
\end{equation}
Its subcritical nonexistence conjecture asks whether the weighted
hyperbola alone rules out positive solutions, without symmetry or
conditions at infinity. The unweighted case \(A=B=0\) is the classical
Lane--Emden system.

The scalar subcritical Liouville theorem of Gidas and Spruck is a
starting point for this question~\cite{GidasSpruck1981}. For systems,
Mitidieri developed a Rellich identity and early nonexistence
regions~\cite{Mitidieri1993,Mitidieri1996}. Serrin and Zou proved the
unweighted three-dimensional assertion under a growth condition;
Pol\'a\v{c}ik, Quittner, and Souplet removed that condition, and
Souplet settled dimension four~\cite{SerrinZou1996,
PolacikQuittnerSouplet2007,Souplet2009}. Li, Li, and Wei obtained a
further unweighted nonexistence region for \(p,q\ge1\) in every dimension
\(n\ge5\)~\cite[Theorem~1.1]{LiLiWei2025}. Liouville theorems also serve
as inputs to rescaling arguments for a priori bounds and singularity
estimates~\cite{GidasSpruckBounds1981,PolacikQuittnerSouplet2007}.

For the weighted system, Bidaut-V\'eron and Giacomini identified the
critical hyperbola in the regular radial theory under admissible
weights~\cite{BidautVeronGiacomini2010}. Phan formulated the unrestricted
weighted question for \(A,B>-2\) as Conjecture~C and established
bounded, low-dimensional, and conditional cases~\cite{Phan}.
Fazly and Ghoussoub obtained bounded and stability-based results for
nonnegative weights~\cite{FazlyGhoussoub2014}; Li and Zhang subsequently
proved the weighted conjecture in dimension three without a boundedness
assumption~\cite{LiZhang2019}. Results in restricted exponent ranges
extend to dimensions four and five~\cite{Li2022}. More recently,
Huang and Zou stated the all-real-weight, origin-continuous assertion
as Conjecture~B and proved cases with stability, decay, or energy
hypotheses~\cite{HuangZou}. Xu and Luo obtained a higher-order weighted
theorem over the full exponent range for radial solutions subject to
regularity and growth conditions~\cite{XuLuo2026}. The full nonradial
question in higher dimensions is the subject of this paper.

\begin{theorem}\label{thm:main}
Let \(n\ge2\) be an integer, let \(p,q>0\), and let \(A,B\in\R\).
Suppose that
\begin{equation}\label{eq:subcritical}
 \frac{n+A}{p+1}+\frac{n+B}{q+1}>n-2.
\end{equation}
There are no functions
\[
 u,v\in C^2(\R^n\setminus\{0\})\cap C(\R^n),
 \qquad u(x)>0,\quad v(x)>0\quad(x\in\R^n),
\]
satisfying Equation~\eqref{eq:system} at every \(x\ne0\).
\end{theorem}

\begin{corollary}[Unweighted Lane--Emden conjecture]
\label{cor:lane-emden}
For every integer \(n\ge3\) and \(p,q>0\) satisfying
\[
 \frac{1}{p+1}+\frac{1}{q+1}>\frac{n-2}{n},
\]
there is no strictly positive \(C^2(\R^n)\) solution of
\[
 -\Delta u=v^p,\qquad -\Delta v=u^q\qquad\text{in }\R^n.
\]
\end{corollary}
